\chapter{Módulos y funcionalidades del sistema}

\section{Portal Público e Interfaz Ciudadana}
El sistema cuenta con un portal institucional que actúa como la puerta de entrada principal para los usuarios. Este módulo no solo centraliza el acceso a la reserva de turnos, sino que funciona como un espacio de divulgación dinámica.

\begin{itemize}
    \item \textbf{Visibilidad de Proyectos e Investigaciones:} Se integra un módulo de publicación donde las dependencias pueden cargar de forma dinámica los proyectos de extensión y las investigaciones en curso, permitiendo a la comunidad conocer las actividades antes de reservar un turno.
\end{itemize}

\begin{figure}[htbp]
    \centering
    % \includegraphics[width=0.8\textwidth]{portal_publico}
    \caption{Placeholder: Captura del Portal Institucional}
    \label{fig:portal_publico}
\end{figure}
\section{Módulo de Gestión de Dependencias y Trámites}
Este módulo constituye el núcleo administrativo del sistema, permitiendo modelar la estructura orgánica de la UNSa y las actividades del Observatorio.

\begin{itemize}
    \item \textbf{Administración de Dependencias:} Los administradores pueden gestionar la jerarquía universitaria. Al crear una dependencia, se define su nombre, código interno y si es una unidad académica. El sistema utiliza el patrón \textit{Nested Set} para permitir que una facultad contenga múltiples institutos o laboratorios.
    \item \textbf{Configuración de Mesas Habilitadas:} A través del \texttt{MesasHabilitadasController}, se gestiona la capacidad operativa de cada oficina. Una "mesa" representa un puesto de atención o un guía disponible. Si el Observatorio cuenta con 3 guías para un horario determinado, se habilitan 3 mesas, lo que limita automáticamente las reservas concurrentes a ese número.
    \item \textbf{Catálogo de Trámites Académicos:} Permite definir actividades con descripciones detalladas, requisitos y estados de activación. Cada trámite se vincula a una o más dependencias, permitiendo una gestión descentralizada.
\end{itemize}

\section{Módulo de Reserva de Turnos}
Diseñado para la simplicidad, este módulo guía al usuario en un flujo interactivo procesado por el \texttt{frontend/TramitesController}.

\begin{itemize}
    \item \textbf{Flujo de Selección:} El solicitante selecciona la Dependencia (ej. Observatorio) y el Trámite (ej. Visita Guiada). El sistema carga dinámicamente los días disponibles.
    \item \textbf{Modalidades de Reserva:} El sistema soporta turnos individuales, grupales, institucionales y mixtos. En modalidades que implican contingentes, se habilita la carga pública de integrantes, donde se solicita información de la institución y una lista de los asistentes.
    \item \textbf{Validación Dinámica:} Antes de mostrar un horario como "Libre", el backend verifica en tiempo real la cantidad de reservas existentes contra las mesas habilitadas. 
    \item \textbf{Comprobante y Notificación:} Una vez confirmada la reserva, se genera un archivo PDF único mediante \texttt{Snappy}. Simultáneamente, se utiliza el sistema de \textbf{Mailables} de Laravel para enviar una confirmación automática.
\end{itemize}

\begin{figure}[htbp]
    \centering
    % \includegraphics[width=0.8\textwidth]{form_reserva_modalidades}
    \caption{Placeholder: Captura del Formulario de Reserva (Modalidades y carga de integrantes)}
    \label{fig:form_reserva_modalidades}
\end{figure}

\section{Módulo de Autogestión de Usuarios ("Mis Turnos")}
Con el objetivo de descentralizar la carga operativa, se implementó un panel personal para que los ciudadanos gestionen sus citas.

\begin{itemize}
    \item \textbf{Panel Mis Turnos:} Permite al usuario autenticado visualizar su historial de reservas y sus turnos próximos.
    \item \textbf{Cancelación y Liberación de Cupos:} El usuario puede cancelar una reserva de forma anticipada especificando un motivo. Esta acción libera instantáneamente el cupo en el sistema, poniéndolo a disposición de otros ciudadanos, y dispara una notificación por correo electrónico.
\end{itemize}

\begin{figure}[htbp]
    \centering
    % \includegraphics[width=0.8\textwidth]{mis_turnos}
    \caption{Placeholder: Captura de pantalla del panel 'Mis Turnos' y proceso de cancelación}
    \label{fig:mis_turnos}
\end{figure}

\section{Módulo de Panel de Operador y Control Diario}
Este módulo está optimizado para el personal que atiende en el Observatorio, contando con un menú lateral reorganizado para facilitar accesos rápidos.

\begin{itemize}
    \item \textbf{Dashboard y Métricas:} Al ingresar, el operador visualiza un panel de control con métricas generales de atención y un resumen de las próximas reservas inmediatas.
    \item \textbf{Gestión de Asistencia y Asignación Manual:} El operador puede marcar el estado del turno (Atendido, Pendiente, Ausente). Adicionalmente, el administrador cuenta con la capacidad de realizar una \textbf{asignación manual de turnos} para casos excepcionales.
    \item \textbf{Cancelación Masiva por Imprevistos:} Ante problemas técnicos o climáticos (frecuentes en observaciones astronómicas), el administrador puede cancelar un bloque de turnos de manera masiva, registrando el motivo y notificando instantáneamente a todos los afectados.
    \item \textbf{Sistema de Pases y Derivaciones:} Permite derivar un ciudadano de una dependencia a otra, manteniendo la trazabilidad.
\end{itemize}

\begin{figure}[htbp]
    \centering
    % \includegraphics[width=0.8\textwidth]{dashboard_admin}
    \caption{Placeholder: Captura del Dashboard Administrativo con métricas}
    \label{fig:dashboard_admin}
\end{figure}

\begin{figure}[htbp]
    \centering
    % \includegraphics[width=0.8\textwidth]{cancelacion_masiva}
    \caption{Placeholder: Captura de la interfaz de Cancelación Masiva o Asignación Manual}
    \label{fig:cancelacion_masiva}
\end{figure}

\section{Módulo de Reportes y Estadísticas}
Para facilitar la toma de decisiones institucionales, el sistema provee herramientas de exportación de datos:
\begin{itemize}
    \item \textbf{Exportación a Excel:} Todos los listados de turnos, feriados y usuarios pueden ser exportados a formato \texttt{.xlsx} para su análisis en herramientas externas.
    \item \textbf{Reportes de Gestión:} Generación de resúmenes de atención por dependencia, permitiendo al Observatorio justificar su carga de trabajo y afluencia de público ante las autoridades universitarias.
\end{itemize}
